\documentclass[a4paper,12pt]{letter}

\usepackage{amsmath}
\usepackage{amssymb}
\usepackage{nopageno}
\usepackage{amsmath}
\usepackage{enumitem}
% \usepackage{mathtools}
\usepackage{eqnarray,amsmath,siunitx}
\usepackage[a4paper, total={7.5in, 11.25in}]{geometry}
%\usepackage[margin=2.5cm]{geometry}
\usepackage[utf8]{inputenc}
\usepackage[T1]{fontenc}
\usepackage{lmodern}
%\usepackage{amsmath, amssymb}
\usepackage{setspace}
\usepackage{hyperref}
\setstretch{1.2}

\begin{document}

\begin{center}
(KTXFI2EBNF) Physics II.  Homework Solutions\\ \smallskip
October~27,~2025
 \end{center}


\textbf{Constants used}
\begin{eqnarray*}
h &=& 6.62607015\cdot10^{-34}\ \mathrm{J\cdot s} \\
c &=& 2.99792458\cdot10^{8}\ \mathrm{m/s} \\
e &=& 1.602176634\cdot10^{-19}\ \mathrm{C} \\
m_e &=& 9.1\cdot10^{-31}\ \mathrm{kg}\quad(\text{value given in the homework})
\end{eqnarray*}

\textbf{Problems and solutions}

\begin{enumerate}
    \item An electron has a kinetic energy of 2.5\,eV.
    \begin{enumerate}
    \item What is its velocity?\\

      Convert energy to joules $E_{kin}=2.5\,\text{eV}=2.5\cdot e$ and use $E_{kin}=\tfrac{1}{2}m_{e}v^{2}$.
      $$v = \sqrt{\frac{2E_{kin}}{m_e}} \approx 9.382520\cdot 10^{5}\,\text{m/s}$$
      
    \item What is its momentum?

    $$p=m_ev\approx 8.538093\cdot 10^{-25}\,\text{kg m/s}$$
      
  \item What is the wavelength of the associated de Broglie wave?
    
    $$\lambda=h/p\approx 7.760597\cdot10^{-10}\,\text{m}$$
\end{enumerate}

 \item A metal surface is illuminated with light of wavelength $2.5 \cdot 10^{-7}\,\text{m}$.  What is the velocity of the emitted electrons if the photoelectric effect starts at light with a wavelength of $2.77 \cdot 10^{-7}\,\text{m}$? $(m_{e} = 9.1 \cdot 10^{-31}\,\text{kg})$\\

 The maximum kinetic energy is $E_{kin}=hc/\lambda - hc/\lambda_{0}$ (if positive), hence
\begin{eqnarray*}
E_{kin} &\approx& 7.744987\cdot 10^{-20}\,\text{J},\\
v &=& \sqrt{\frac{2E_{kin}}{m_{e}}}\approx 4.125767\cdot 10^{5}\,\text{m/s}.
\end{eqnarray*}


\item A photon ejects an electron with a maximum kinetic energy of 0.64\,eV from a metal for which the work function is 3.74\,eV.
\begin{enumerate}
\item What is the energy of the photon in electronvolts?

  $$E_\text{ph}=K+\phi=4.38\ \text{eV.}$$
  
\item What is the wavelength of the applied ultraviolet radiation?

$$\lambda=hc/E_\text{ph}\approx 2.830689\cdot 10^{-7}\,\text{m.}$$
\end{enumerate}

\pagebreak
\item In a Compton scattering experiment, X-rays with a wavelength of $0.104\,\text{nm}$ are used. $(m_{e} = 9.1 \cdot 10^{-31}\,\text{kg})$
    \begin{enumerate}
    \item At what scattering angle does the wavelength of the radiation increase by 1\%?

      \smallskip
      The Compton formula $\Delta\lambda=\lambda_{C}(1-\cos\theta)$ with $\lambda_{C}=h/(m_{e}c)$ gives $\theta \approx 55.1237^{\circ}$.
     
        \item At what angle does the wavelength become 0.05\% larger?\\$\theta \approx 11.8774^{\circ}$.
    \end{enumerate}

 \item X-rays are scattered by electrons. $\lambda_0 = 10^{-11}\,\text{m}$. The magnitude of the wavelength change is $2.6 \cdot 10^{-12}\,\text{m}$.
    \begin{enumerate}
    \item What is the scattering angle of the photons?\\

      Using $\Delta\lambda=\lambda_{C}(1-\cos\theta)$, $\theta \approx 94.0417^{\circ}$.
    
    \item By how much did the photon energy change during the process? $(m_e = 9.1 \cdot 10^{-31}\,\text{kg})$\\

Photon energies before and after are $E_{i}=hc/\lambda_{0}$ and $E_{f}=\frac{hc}{\lambda_{0}+\Delta\lambda}$, so the change is $$\Delta E = E_{i}-E_{f} \approx 4.099015\cdot 10^{-15}\,\text{J}=25584.040947\,\text{eV}.$$
    \end{enumerate}

\end{enumerate}

\bigskip

\rightline{Dr.~Gabor~FACSKO}
\rightline{\textit{facsko.gabor@uni-obuda.hu}}

\vfill


\end{document}
